\chapter{الألم: إشارة إنذار}
\chaptermark{الألم}
\label{sec:pain}

\section{إنذار لا قياس}

كل مستشعرات الفصل~\ref{sec:sensors} كانت تقيس العالم: كم من الضوء، وما حدة الصوت، وما الضغط. أما الألم فمبني على نحو آخر. إنه لا يُبلغ ”ماذا هناك“، بل ”خطر، اترك كل شيء“. وهو عند المهندس ليس قناة قياس، بل \emph{إشارة إنذار}: خط ذو أولوية عليا يجب أن يخترق أي عمل تؤديه الآلة في تلك اللحظة.

يختلف نظام الإنذار عن قناة القياس بثلاث خصائص، وكلها في الألم. أولًا، يصعب إخماده بالتعوّد: مستشعرات الألم تتعوّد أضعف بكثير من مستشعرات اللمس، بل تصبح بعد إصابة خفيفة أشد حساسية~\cite{LaMotte1982hyper}؛ ولم يُقَس ضعف الاستجابة بعد تسخين قوي إلا لدى نوع واحد منها~\cite{Treede1995heat}. مستشعر الضوء يصمت على خلفية ثابتة (الشكل~\ref{fig:adaptation})، أما جهاز إنذار يصمت بعد دقيقة فلا فائدة منه. ثانيًا، عتبته عالية: لا تطلق مستشعرات الألم إلا حين يقترب المؤثر من حد الخطر، فللتسخين مثلًا تتراوح عتبات المستشعرات المختلفة من $41^\circ$ إلى $46\text{--}48^\circ$ بحسب النوع~\cite{Treede1995heat, Weidner1999cfib}. ثالثًا، وهذا هو الأهم في هذا الفصل، لا يقرر علوَّ صوت الإنذار المستشعرُ وحده، بل الآلة نفسها أيضًا. فالجرح نفسه يؤلم على نحو مختلف في ظروف مختلفة. فلنرَ من أي قطع مألوفة لنا جُمع ذلك.

مستشعرات الألم عصبونات \nt{GLUT}، كبقية العصبونات الحسية في الفصل~\ref{sec:sensors}. وبعضها يطلق مع الغلوتامات المادة~\foreignlanguage{english}{P} من الملحق~\ref{sec:othermediators}، التي تقوّي النقل ببطء.

\section{سلكان: ألم سريع وألم بطيء}

اضرب إصبعك، وستشعر بالألم مرتين: أولًا ألم قصير حاد، ثم بعد ثانية ألم كليل طويل. هذان سلكان مختلفان. السلك السريع معزول، ويوصل النبضات بسرعة $v$ من $5$ إلى $30$ مترًا في الثانية؛ والبطيء رفيع وعارٍ، $0.5\text{--}2$ متر في الثانية. تقطع الإشارة من القدم نحو متر، وزمن الوصول
\begin{empheq}[box=\keybox]{equation}
t \;=\; \frac{L}{v}
\label{eq:paintime}
\end{empheq}
للسلك السريع عشرات الملّي ثانية، وللبطيء نحو ثانية. يحمل السلكان رسالتين. السريع يقول \emph{أين} و\emph{متى}: نبضاته تصل أبكر وأدق، كما في شفرة التوقيت من الفصل~\ref{sec:code}. والبطيء يقول \emph{ما مدى السوء}، وألمه الكليل الطويل يُبقي الآلة في نمط الإنذار حتى تزول الإصابة.

\section{البوابة في الحبل الشوكي}

أول ما تصل إليه إشارة الألم الحبل الشوكي، وهناك تمر عبر \emph{بوابة}~\cite{MelzackWall1965}؛ وقد وُجدت العصبونات المحددة لهذه البوابة حديثًا نسبيًا~\cite{Duan2014}. على عصبون \nt{GLUT} المخرج، الذي ينقل الألم إلى الدماغ، يلتقي سلك الألم وسلك اللمس. وبجانبه عصبون وسيط مثبِّط، \nt{GABA} أو \nt{GLYC}، يثيره اللمس (الشكل~\ref{fig:gate}). فاللمس يغلق البوابة. وفي نظرية البوابة الأصلية~\cite{MelzackWall1965} يُطفئ الألمُ على العكس هذا الوسيط، فيفتح الألم القوي البوابة؛ وهذا افتراض في النظرية، لم يُبيَّن مباشرة على عصبونات البوابة المكتشفة.

\begin{figure}[t]
\centering
\begin{tikzpicture}[>=Stealth,
  nrn/.style={circle,draw,very thick,minimum size=1.0cm,inner sep=0pt,font=\scriptsize},
  inh/.style={-{Bar[width=7pt]},thick,red!70!black},
  bc/.style={->,thick,densely dashed,orange!85!black}]
\node[nrn,fill=blue!6] (Z) at (6,0) {\nt{GLUT}};
\node[nrn,fill=red!10] (Q) at (3.6,-1.6) {\nt{GABA}};
\draw[->,very thick,red!70!black] (0,0.6) node[left,font=\small,black,align=right] {الألم $\nu$} -- (5.45,0.15);
\draw[->,very thick,blue!60!black] (0,-2.6) node[left,font=\small,black,align=right] {اللمس $\mu$} -- (2.9,-2.0);
\draw[->,thick,blue!60!black] (1.8,-2.35) -- (5.5,-0.35) node[pos=0.8,below right=-2pt,font=\scriptsize] {$w_{z\mu}$};
\draw[inh] (1.6,0.5) -- (3.35,-1.1) node[pos=0.5,left=1pt,font=\scriptsize] {$w_{q\nu}$};
\draw[inh] (Q) -- (Z) node[pos=0.5,above left=-1pt,font=\scriptsize] {$\beta$};
\draw[->,very thick] (Z) -- (8.2,0) node[right,font=\small,align=left] {إلى الدماغ: $z$};
\draw[bc] (6,2.1) node[above,font=\scriptsize,black,align=center] {من الأعلى: \nt{SERO}، \nt{NORA}، الأفيونيات} -- (Z.north) node[pos=0.45,right,font=\scriptsize,black] {الكسب $g$};
\end{tikzpicture}
\caption{بوابة الألم في الحبل الشوكي بحسب النموذج~\eqref{eq:gate}. يتلقى عصبون \nt{GLUT} المخرج الألم $\nu$ وإثارة ضعيفة من اللمس $\mu$. الوسيط المثبِّط (\nt{GABA} أو \nt{GLYC}) يثيره اللمس، ويطفئه الألم بحسب افتراض نظرية البوابة. ومن الأعلى، عبر قنوات البث، يأتي الكسب $g$.}
\label{fig:gate}
\end{figure}

لنكتب ذلك للترددات، كما في الفصل~\ref{sec:gaba}. تردد الوسيط $q$ وتردد المخرج $z$ هما
\begin{empheq}[box=\keybox]{equation}
q \;=\; \bigl[w_{q\mu}\,\mu + q_0 - w_{q\nu}\,\nu\bigr]_+ ,
\qquad
z \;=\; \bigl[\,g\,(\nu + w_{z\mu}\,\mu) - \beta\,q\,\bigr]_+ ,
\label{eq:gate}
\end{empheq}
حيث $\nu$ دخل الألم، و$\mu$ دخل اللمس، و$w$ أوزان الوصلات (الدليل الأول: إلى مَن، والثاني: من مَن)، و$\beta$ شدة التثبيط، و$q_0$ النشاط الخلفي الذاتي للوسيط، و$g$ الكسب الذي يُحدَّد من الأعلى (عنه أدناه). هذا هو المنع~\eqref{eq:veto} من الفصل~\ref{sec:gaba}: ”الألم، ولكن ليس مع اللمس“، مع تعديل واحد مأخوذ من نظرية البوابة افتراضًا: الألم القوي يطفئ بنفسه العصبون المانع.

حسبنا النموذج عند $w_{q\mu}=1$، $q_0=0.2$، $w_{q\nu}=0.5$، $w_{z\mu}=0.3$، $\beta=1.5$ (الشكل~\ref{fig:gatecurves}). بلا لمس يبدأ الألم بالمرور من $\nu\approx0.18$. ومع اللمس ترتفع العتبة إلى $0.86$، وعند $\nu=1$ يهبط المخرج أربع مرات، من $1$ إلى $0.25$. أما مع الألم القوي، $\nu=2$، فلا يغيّر اللمس شيئًا: البوابة مفتوحة على مصراعيها. وهكذا في الحياة: فرك الكدمة يساعد، أما في الإصابة الخطيرة فلا. واللمس نفسه، بلا ألم، لا يمر عبر البوابة: الإنذار لا ينطلق من اللمس.

\begin{figure}[t]
\centering
\begin{tikzpicture}[>=Stealth,x=5.2cm,y=1.35cm]
\draw[->,gray!70!black] (0,0) -- (2.12,0) node[right,font=\small,black] {الألم $\nu$};
\draw[->,gray!70!black] (0,0) -- (0,4.3) node[left,font=\small,black] {$z$};
\foreach \t in {0,0.5,1,1.5,2}{\draw[gray!60!black] (\t,-0.05) -- (\t,0.05); \node[below,font=\scriptsize] at (\t,0) {\t};}
\foreach \f in {1,2,3,4}{\draw[gray!60!black] (-0.01,\f) -- (0.01,\f); \node[left,font=\scriptsize] at (0,\f) {\f};}
\draw[very thick,red!70!black] plot file {figures/pain_gate_notouch.dat};
\draw[very thick,blue!60!black] plot file {figures/pain_gate_touch.dat};
\draw[very thick,orange!85!black,densely dashed] plot file {figures/pain_gate_sens.dat};
\draw[very thick,red!70!black] (0.08,3.9) -- (0.2,3.9) node[right,font=\scriptsize,black] {بلا لمس};
\draw[very thick,blue!60!black] (0.08,3.45) -- (0.2,3.45) node[right,font=\scriptsize,black] {مع اللمس};
\draw[very thick,orange!85!black,densely dashed] (0.08,3.0) -- (0.2,3.0) node[right,font=\scriptsize,black] {بعد التحسّس، $g=2$};
\end{tikzpicture}
\caption{مخرج البوابة~\eqref{eq:gate} بحسب شدة الألم. اللمس يرفع العتبة ويخفت الألم الضعيف والمتوسط، لا القوي. والتحسّس (الخط المتقطع) يضاعف الكسب: الألم نفسه يُنقل بصوت أعلى مرتين، وتنخفض العتبة.}
\label{fig:gatecurves}
\end{figure}

\section{مقبض الصوت من الأعلى}

لا تُدار البوابة من الأسفل وحده. فمن جذع الدماغ تنزل إليها مسارات هابطة تغيّر الكسب $g$ في~\eqref{eq:gate}: \nt{SERO} و\nt{NORA} والببتيدات الأفيونية من الملحق~\ref{sec:othermediators}~\cite{Fields2004}. هذه هي قنوات البث نفسها التي في الفصول من~\ref{sec:serocomputer} إلى~\ref{sec:acet}، إلا أن مقبضها هنا علوّ صوت الإنذار. وهي تستطيع إدارته في الاتجاهين: الدارات الهابطة نفسها قادرة على إخفات الألم وعلى تضخيمه.

لماذا تخفض الآلة الإنذار؟ لأن الإنذار مقاطعة، والمقاطعة ليست مناسبة دائمًا. الحيوان الهارب من مفترس يجب ألا يتوقف بسبب قائمة جريحة: الركض أولًا، ثم لعق الجرح. وتوقع الراحة أيضًا يخفت الألم، ويزول جزء من هذا الأثر إذا حُجبت المستقبِلات الأفيونية~\cite{Levine1978}: فالتوقع المحسوب في الدماغ ينزل عبر القناة الأفيونية إلى البوابة. هذه الصورة بحد ذاتها مقيسة؛ أما كيف يقرر الدماغ بالضبط كم يجب أن يكون علو الصوت ففرضية.

\section{التحسّس: إنذار يتعلم}

بعد الإصابة تصبح عصبونات المخرج في البوابة أشد حساسية: الدخل نفسه يستدعي مخرجًا أكبر، والعتبة تنخفض~\cite{Woolf1983}. في النموذج~\eqref{eq:gate} هذا ازدياد في الكسب $g$، وأبسط وصف له دارة \foreignlanguage{english}{RC} بطيئة يشحنها الألم نفسه:
\begin{empheq}[box=\keybox]{equation}
\tau_s\,\frac{\dd g}{\dd t} \;=\; -(g-1) + k_s\,z ,
\label{eq:sensitization}
\end{empheq}
حيث $\tau_s$ الزمن الذي تعود فيه الحساسية إلى طبيعتها، وهو أطول من مدة الألم نفسه، لأن التحسّس يدوم حتى بعد توقف المنبّه المؤذي~\cite{Woolf1983}، و$k_s$ قوة الشحن. ما دام النسيج مصابًا يُبقي مخرج البوابة $g$ فوق الواحد، فينتقل كل ما يحدث قرب الجرح بصوت أعلى (الخط المتقطع في الشكل~\ref{fig:gatecurves}: عند $g=2$ تنخفض العتبة إلى $0.11$، ويتضاعف المخرج). هذا ضبط معقول لجهاز الإنذار: ما دامت الإصابة لم تلتئم، فالاحتياط أفضل.

هذه عند المهندس تغذية راجعة موجبة بتسريب بطيء: الألم يرفع الكسب، والكسب يرفع الألم. ما دامت الحلقة ضعيفة يعود $g$ إلى طبيعته حين يلتئم الجرح. أما إذا كانت الحلقة قوية، فقد يعلق جهاز الإنذار في حالة التشغيل حتى بعد زوال الإصابة، كالمجموعة ذات التغذية الراجعة القوية في الشكل~\ref{fig:bistable}. النموذج~\eqref{eq:sensitization} تبسيط؛ أما وجود التحسّس المركزي فمقيس.

\section{الألم معلّمًا ومقاطعةً}

للألم في الآلة عملان إلى جانب الإنذار.

\emph{المعلّم.} الألم مكافأة سالبة: يدخل في صيغة الخطأ~\eqref{eq:ctd} بوصفه $r<0$. وكل ما قيل في الفصل~\ref{sec:dofa} يصح هنا أيضًا: نذير الألم يبدأ باستدعاء الخطأ مسبقًا، فتتعلم الشبكة تجنب ما يقود إلى الألم. والتجارب على البشر تُظهر أن نشاط الدماغ في التعلم من الألم يتبع خطأ الفروق الزمنية بالذات~\cite{Seymour2004}، وأن بعض عصبونات \nt{DOPA} تستجيب للأحداث المزعجة أيضًا (الفصل~\ref{sec:dofaalt}).

\emph{المقاطعة.} الألم يصرف الانتباه عن المهمة الجارية، وهذه وظيفته لا أثر جانبي له~\cite{EcclestonCrombez1999}. وفي الفصل~\ref{sec:nora} نوقش دور ”المقاطعة“ نفسه للرشقة الطورية من \nt{NORA}. أما أسرع استجابة فلا تصل إلى الدماغ أصلًا: سحب اليد عن الساخن يُغلق مباشرة في الحبل الشوكي بزوج العضلات نفسه والمنع نفسه للعضلة المضادة كما في الشكل~\ref{fig:antagonists}.

\begin{keyblock}
\begin{proposition}[إشارة الإنذار]
\label{prop:pain}
الألم ليس قياسًا، بل إشارة إنذار ذات أولوية عليا. يتعوّد أضعف بكثير من قنوات القياس، ويسير في سلكين (السريع: ”أين“، والبطيء: ”ما مدى السوء“)، ويمر عبر بوابة يغلقها اللمس (ويفتحها الألم القوي بحسب افتراض نظرية البوابة)، وتحدد علوَّ صوته من الأعلى قنواتُ البث. هذه الأجهزة مقيسة، عدا فتح البوابة بالألم؛ والنموذجان البسيطان~\eqref{eq:gate} و\eqref{eq:sensitization} تبسيطان.
\end{proposition}
\end{keyblock}

\section{الخلاصة}

\begin{table}[t]
\centering
\footnotesize
\setlength{\tabcolsep}{4pt}
\renewcommand{\arraystretch}{1.15}
\begin{tabular}{>{\raggedright}p{3.0cm}>{\raggedright}p{4.9cm}>{\raggedright\arraybackslash}p{4.9cm}}
\toprule
ماذا يفعل الألم & كيف & ما هو معروف \\
\midrule
يطلق الإنذار & عتبة عالية، وتعوّد أضعف من اللمس & العتبة مقيسة؛ ومقارنة التعوّد تقدير \\
يُبلغ ”أين“ و”ما مدى السوء“ & سلكان بسرعتين مختلفتين~\eqref{eq:paintime} & مقيس \\
يمر عبر بوابة & منع عبر \nt{GABA}/\nt{GLYC}~\eqref{eq:gate} & البوابة مقيسة؛ وفتحها بالألم افتراض النظرية؛ والنموذج تبسيط \\
يغيّر علو الصوت & \nt{SERO} و\nt{NORA} والأفيونيات الهابطة & مقيس؛ وقاعدة اختيار علو الصوت فرضية \\
يتعلم بعد الإصابة & ازدياد الكسب~\eqref{eq:sensitization} & التحسّس مقيس؛ والنموذج تبسيط \\
يعلّم التجنب & مكافأة سالبة في~\eqref{eq:ctd} & التشابه مع خطأ الفروق الزمنية مقيس \\
يقاطع العمل & الاستيلاء على الانتباه؛ منعكس السحب & مقيس \\
\bottomrule
\end{tabular}
\caption{الألم إشارةَ إنذار.}
\label{tab:pain}
\end{table}

الألم مجمَّع من قطع رأيناها من قبل (الجدول~\ref{tab:pain}): مستشعرات \nt{GLUT}، وسلكان، ومنع عبر وسيط مثبِّط، ومقبض صوت يُدار بالبث، ودارة \foreignlanguage{english}{RC} بطيئة، وخطأ التنبؤ، والمقاطعة. والجديد فيه أمر واحد: إنه المدخل الوحيد للآلة الذي تحدد الآلة نفسها علوّ صوته، فهو جهاز إنذار يستطيع صاحبه أن يخفته وأن يعيد ضبطه.

\section{تمارين}

\begin{exercise}[\lvlA]
\label{ex:14:1}
\emph{سلكان.} تأتي الإشارة من القدم، $L=1\,\text{m}$. (أ)~تسير الدفعة في حزمة من الأسلاك من نوع واحد، تملأ سرعاتها المدى~\eqref{eq:paintime}. بكم تتمدد حين تصل إلى الحبل الشوكي، للنوع السريع وللبطيء؟ (ب)~موجتا الألم في سلكين سرعتاهما $15$ و$1\,\text{m/s}$ يمكن تمييزهما إذا كان الفرق $0.1\,\text{s}$ فأكثر. من أي مسافة تندمجان؟
\end{exercise}

\begin{exercise}[\lvlB]
\label{ex:14:2}
\emph{متى يؤلم اللمس.} البوابة~\eqref{eq:gate} بمعاملات النص، والكسب $g$ يزداد مع التحسّس. حتى أي $g$ لا يمر اللمس بلا ألم عبر البوابة مهما كانت شدة $\mu$؟ وعند أي $g$ يبدأ اللمس $\mu=1$ بالمرور؟
\end{exercise}

\begin{exercise}[\lvlB]
\label{ex:14:3}
\emph{استقرار التحسّس.} المداخل ثابتة، ومخرج البوابة خطي في $g$: $z=gA-C$، $A=\nu+w_{z\mu}\mu$، $C=\beta q\ge0$. (أ)~أوجد التوازن $g^*$ للنموذج~\eqref{eq:sensitization} وشرط استقراره. (ب)~عند $k_s=0.5$ أوجد شدة الألم التي ينفلت النموذج فوقها، بلا لمس ومع اللمس $\mu=1$.
\end{exercise}

\begin{exercise}[\lvlA]
\label{ex:14:4}
\emph{نذير الألم.} إشارة في اللحظة $0$ تتنبأ بألم في اللحظة $T=2\,\text{s}$، وفي~\eqref{eq:ctd} $r(t)=-P\,\delta(t-T)$، $\tau_\gamma=2\,\text{s}$. (أ)~أوجد مساحة دفقة $\delta$ لحظة الإشارة. (ب)~فعل في اللحظة $t_e=1\,\text{s}$ يلغي الألم. ما دفقة $\delta$ في هذه اللحظة، وماذا ستعلّم الشبكة؟
\end{exercise}

\begin{exercise}[\lvlB]
\label{ex:14:5}
\emph{إنذار ينغلق على نفسه.} أضف إلى البوابة من \foreignlanguage{english}{\texttt{models/\allowbreak pain\_gate.py}} ديناميكا~\eqref{eq:sensitization} وتشبعًا $z\le4$. اللمس $\mu=1$ موجود دائمًا، والإصابة $\nu=2$ تدوم زمنًا $T$، ثم $\nu=0$. أوجد بالنمذجة أصغر $T$ (بوحدات $\tau_s$) يبقى الإنذار بعده مشغَّلًا، عند $k_s=4$، $4.6$، $5$، $6$، $8$.
\end{exercise}

\begin{exercise}[\lvlB]
\label{ex:14:6}
\emph{تصميم البوابة.} عند $\beta=1.5$، $w_{z\mu}=0.3$، $g=1$ اختر في~\eqref{eq:gate} القيم $q_0$ و$w_{q\nu}$ و$w_{q\mu}$ بحيث تكون عتبة الألم $0.2$ بلا لمس و$1.0$ مع اللمس $\mu=1$، ويكفّ اللمس عن إخفات الألم عند $\nu_\times=2.5$. حتى أي كسب $g$ لا يمر اللمس بلا ألم عبر البوابة؟
\end{exercise}

\begin{exercise}[\lvlC]
\label{ex:14:7}
\emph{مقبض الصوت.} العمل يجلب قيمة $V$ في وحدة الزمن، والعمل مع إصابة $\nu$ يكلف $c\nu$، فتكون المقاطعة مجدية عند $\nu>V/c$؛ والألم يقاطع عند $z>\theta=0.5$. (أ)~أي كسب $g^*$ للبوابة~\eqref{eq:gate} بلا لمس يعطي هذه العتبة عند $V/c=0.25$، $0.5$، $2$؟ (ب)~من أي إشارات الكتاب يمكن للآلة أن تقدّر $V$ و$c$؟
\end{exercise}
